\section{Introduction} \label{sec:introduction}
In a world of increasing amounts of electric 
vehicles, data centres, and 
other large-scale electricity-hungry industries, electricity grids are 
beginning to feel the strain~\cite{ppm-gridImpact}, contributing to an increase 
in electricity prices~\cite{gridStrainIncreasingGraph}. When faced with such 
increased prices, it is beneficial, from a monetary point of view, to schedule 
electricity intensive tasks during periods where prices are low. 
However, to do this, one must know when such low-price periods occur. In 
Denmark, 
the price of electricity is, in simple terms, fixed for the next 24-hour period 
at noon daily\footnote{This is a simplification for ease of reading. For actual 
	energy trading details one may visit 
	\cite{nordpoolTrading}. It is also further 
	explained in \autoref{sec:prereq}}. 
This is often sufficient for household electricity use, as one can typically 
schedule dishwashers, washing machines, and electric vehicle charging to happen 
during the lowest point of that period. This can be done either through direct 
scheduling on individual appliances or smart devices, or through smart home 
management systems such as Home Assistant~\cite{homeAssistant}. Industrial 
needs for 
energy can involve larger and more time-sensitive electricity 
demand~\cite{industrialHouseholdEnergyConsumption}, and may 
therefore be more expensive. As an example, one might imagine a 
near-future where fleets of electric trucks need their batteries charged 
before delivery. In a scenario like this, it may be beneficial for such 
companies, to be able to forecast electricity prices further into the future, 
and plan their electricity consumption accordingly, to minimize charging costs. 
\\

When attempting to predict electricity prices, many different aspects can be 
considered~\cite{forecastInput}. Among these, the ones of particular interest 
to this paper are: 
\begin{itemize}
	\item Weather
	\item Electricity production details
	\item Electricity consumption details
	\item Historical- or T-minus 
	data\footnote{Data from an earlier point in 
		time before the current time, e.g. 24 hours 
		before.}
	\item Time of year/month/week/day
	\item Electricity exchange with foreign countries
\end{itemize}

One factor that influences electricity prices is the weather, 
due to its direct impact on both electricity 
production~\cite{productionPatterns} and 
consumption~\cite{consumptionPatterns}. 
In a country that has invested extensively in wind-energy, such as Denmark, the 
weather directly impacts electricity production, as more that $80\%$ is 
produced by renewable 
sources~\cite{energiStyrelsenEnergiStatestik2022}. Steady winds, and to a 
lesser extent, clear skies therefore affect energy production quite heavily in 
Denmark, as necessities for efficient wind and solar energy production. 
When it comes to consumption, human consumption patterns also change depending 
on the weather~\cite{consumptionPatterns}. When it 
is cold, people turn on the 
heat to warm themselves, 
and when it is hot, they turn on air conditioning, or fans, to stay cool. 
Further, when it is cold, such as in the winter months, it is typically also 
darker, necessitating the need to use lights more frequently. 

These factors suggest that one can use the weather to predict electricity 
\gls{pc} patterns, which implies that weather data may carry 
useful information for electricity price forecasting. A substantial body of 
prior work has examined short-term electricity price forecasting using 
statistical, machine-learning, and deep-learning 
approaches~\cite{energyPricePredictionMLNeighbors,energyPricePredictionUK,energyPricePredictionNovelDeepLearning},
showing various results in different forecasting scenarios. However, many of 
the reported results are obtained in specific market settings or with highly 
specialized forecasting architectures, such as in 
\cite{energyPricePredictionNovelDeepLearning}. Much of this prior work shares 
common limitations, such as results typically being reported for a single model 
trained on a fixed feature set. This may obscure the answer to what data 
actually impact the performance of models in general. Someone wishing to deploy 
a forecasting system is therefore left without clear guidance on whether 
investing in comprehensive data collection is worthwhile, or whether a simpler 
feature set would perform comparably. Further, while related studies exist, the 
Danish setting combines several factors that are not jointly emphasized in most 
of the literature reviewed here, such as its heavy investment into 
wind-energy~\cite{energiStyrelsenEnergiStatestik2022}, its division of the 
energy grid into two distinct bidding zones, and its heavy energy-exchange with 
its neighbouring countries. \\

This paper addresses these gaps through a large-scale ablation study, 
systematically varying the feature composition, forecasting horizon, and 
prediction target across 1.052 trained models. Rather than reporting the best 
achievable performance under ideal conditions, the goal is to answer a 
practically motivated question: how much data do you actually need? This is not 
merely an academic exercise, but a question of practicality. An energy trading 
company deploying a real-time forecasting system must weigh the cost of 
acquiring and maintaining each data source against the marginal accuracy it may 
provide. Understanding which feature categories contribute meaningfully, and at 
which horizons they stop being useful, has direct commercial value.\\

To investigate this dynamic across a broad spectrum of architectures, eight 
different types of general-purpose machine learning algorithms and frameworks 
are evaluated: 
\begin{multicols}{2}
	\begin{enumerate}
		\item RandomForest
		\item XGBoost
		\item LightGBM
		\item CatBoost
		\item LSTM
		\item GRU
		\item Transformer
		\item AutoGluon
	\end{enumerate}
\end{multicols}
Model training will be based on weather data gathered from the \gls{dmi}, as 
well as electricity \gls{pc} data from \gls{en}. 
To comprehensively evaluate the chosen forecasting 
approaches, the experiments are structured around 
three distinct prediction settings. First, the 
models are trained on various data 
constellations, ranging from purely 
weather-related data, to a more comprehensive 
dataset of weather, grid, and historical 
price data. This is done to isolate the impact of 
each specific data category. Second, the 
forecasting horizon is expanded sequentially to 
evaluate how the accuracy of each algorithm 
deteriorates as it predicts further into the 
future. Finally, the approaches are trained using 
two distinct prediction targets: predicting the 
absolute electricity price and predicting the 
price delta. For the purpose of this paper, the 
price delta refers to the change in price from the 
current hour to the forecasted hour. This 
dual-target approach is motivated by a desire to 
anticipate model overfitting and to investigate 
whether the algorithms are more adept at capturing 
relative price fluctuations rather than absolute 
prices. The overall purpose of this ablation study 
is to uncover how different machine learning 
frameworks react when given identical training 
conditions under these varying settings.\\

Producing a meaningful ablation study in this setting required substantial 
preprocessing effort. The raw data sources used for this paper spans nearly 11 
years, and consists of station-level weather measurements, energy production 
and consumption records, and price data across two format changes. This data 
was not designed for machine learning consumption. Constructing a consistent, 
temporally aligned training dataset required geographic filtering of weather 
stations, cross-year column homogenisation, resolution standardisation, and 
imputation of structurally missing data. This preprocessing pipeline is a 
necessary foundation for the validity of the experimental results, and is 
described in detail in \autoref{sec:data_preprocessing} to enable 
reproducibility\hltodo{.}{Will likely need to shorten introduction, though i am 
not entirely sure how, as much of this motivation is necessary.}\\

Related to this, one company has 
actually already utilised weather data to 
predict electricity production and consumption in Denmark. This company is 
called Midas Energy~\cite{midasEnergy}. They do this with the purpose of 
supplying their customers, energy trading companies, with this information, in 
order to enable them to make decisions on a stronger foundation. They provide 
short term energy production and consumption forecasts, so that their 
customers may buy and sell electricity with the ability to consider the future 
to a greater degree than they otherwise would be able to~\cite{midasEnergy}. 
Notably, Midas Energy provides production and consumption forecasts, not price 
forecasts, but is interested in expanding into such areas, which is why they 
have acted as partners during the development process of this paper. Their 
contributions have been mostly related to data acquisition, performance metric 
evaluations, and benchmarking, and their contributions will be discussed 
further when relevant.